\section{Selection on a common data domain}
\label{h:alternatives}

The two applications use the same analytic selection step: several
continuation measures are supported on one data domain, and a measurable
route chooses among them. The strong density estimate of
Theorem~\ref{h:dv:finite-density} is sufficient, but a bound for
measurable tests already gives the required \(N^{1/p}\) cost.

We state this common step before applying the projection theorem to a
lower-dimensional cube prior in Section~\ref{h:cube:section}.
The later sections examine additional exact-label interfaces and
alternative coefficient normalizations; they do not assert a stronger
main theorem.

\begin{lemma}[Selection on a common data domain]\label{h:alt:selection}
Let \((E,\mathcal E,\lambda)\) be a measure space, let
\(\mathcal D\in\mathcal E\) have finite measure, let \(p>1\), let
\(U\ge0\), and let \(N\ge1\) be an integer. Suppose finite nonnegative measures
\(\eta_1,\ldots,\eta_N\) are supported on \(\mathcal D\), and that
for every measurable \(0\le g\le1\),
\begin{equation}\label{h:alt:test}
 \int g\,d\eta_j
 \le U\left(\int_{\mathcal D}g\,d\lambda\right)^{1-1/p}
 \qquad(1\le j\le N).
\end{equation}
For measurable routes \(g_j\ge0\) with \(\sum_jg_j\le1\),
\begin{equation}\label{h:alt:selection-bound}
 \sum_{j=1}^N\int g_j\,d\eta_j
 \le N^{1/p}U\,\lambda(\mathcal D)^{1-1/p}.
\end{equation}
A sufficient condition for \eqref{h:alt:test} is
\(\eta_j=f_j\lambda\), where \(f_j\ge0\) vanishes outside
\(\mathcal D\) and \(\|f_j\|_{L^p(\lambda)}\le U\).

In particular, if \(r>0\), \(m,\beta\in\R\), and positive constants
\(H,V,L\) satisfy
\[
 U\le HLr^{\beta-m(1-1/p)},\qquad
 \lambda(\mathcal D)\le Vr^m,
\]
then the right side of \eqref{h:alt:selection-bound} is at most
\begin{equation}\label{h:alt:scale-cancellation}
 N^{1/p}H V^{1-1/p}Lr^\beta.
\end{equation}
\end{lemma}

\begin{proof}
Each \(g_j\) lies in \([0,1]\). Put
\(b_j=\int_{\mathcal D}g_j\,d\lambda\). Equation
\eqref{h:alt:test} and H\"older on the finite index set give
\[
 \sum_j\int g_j\,d\eta_j
 \le U\sum_jb_j^{1-1/p}
 \le U N^{1/p}\left(\sum_jb_j\right)^{1-1/p}.
\]
Since \(\sum_jg_j\le1\), we have
\(\sum_jb_j\le\lambda(\mathcal D)\), proving
\eqref{h:alt:selection-bound}. Under the density condition,
H\"older gives
\[
 \int g\,d\eta_j
 \le U\left(\int_{\mathcal D}g^{p/(p-1)}\,d\lambda\right)^{1-1/p}
 \le U\left(\int_{\mathcal D}g\,d\lambda\right)^{1-1/p},
\]
because \(0\le g\le1\). Finally the powers of \(r\) in the two
assumed bounds add to \(\beta\), giving
\eqref{h:alt:scale-cancellation}.
\end{proof}

The test formulation accommodates a direct open-set limit even when
one chooses not to infer a strong \(L^p\) density from that limit.
The density formulation applies to the other passages. In both cases,
the reference measure and the common support must be the ones
specified by the corresponding projection estimate.
